\section{Legal Posture Mapping}
\label{sec:legal-posture}

The insensitivity finding in Section~\ref{sec:results} means that the choice
of percentile $p$ is primarily a \emph{legal and normative} choice rather
than an empirical one.
Different values of $p$ correspond to different legal doctrines and different
rhetorical postures before a court.
This section maps three postures and recommends one for statutory adoption.

\subsection{Posture 1: Compactness Doctrine ($p = 0.0$)}

Setting $p = 0.0$ causes \ps{} to behave as \cs{}: it returns the
minimum-edge-cut plan across all $T$ seeds.
This is the plan that is more compact than every other plan in the bisection
family.

The legal foundation is \karcher{}~\citep{karcher1983} and the compactness
tradition.
\karcher{} requires that any deviation from strict population equality be
justified by a legitimate state objective.
By analogy, compactness doctrine requires that any departure from maximum
compactness be justified by a legitimate purpose.
The $p = 0.0$ plan satisfies this standard trivially: there is no plan in
the bisection family that is more compact.

\begin{statute}
\textbf{Compactness Posture.}
The Districting Administrator shall produce the plan with the smallest
normalised edge-cut fraction among the $T$-plan bisection distribution.
This plan is the most compact consistent with population balance and
contiguity within the \ar{} bisection family.
No neutral party can produce a more compact plan within this family.
\end{statute}

The $p = 0.0$ posture has a powerful adversarial property.
Under \karcher{}, any challenger must produce an alternative plan that
is at least as compact.
For states where the $p = 0.0$ plan sits at the 0.1--0.2nd percentile
of the \recom{} ensemble (WI, GA, PA per G.1~\citep{g1paper}), this means
the challenger must produce a plan more compact than 99.8\% of all valid
plans.
This is an extraordinarily high bar.

The potential vulnerability is the over-optimising objection: the algorithm
was specifically designed to find the extremum.
However, this objection fails on its own terms.
Compactness is not a zero-sum quantity: maximising compactness does not
require minimising anything else.
The $p = 0.0$ plan minimises boundary crossings, which directly serves the
interest of communities that prefer not to be split across districts.

\subsection{Posture 2: Representativeness Doctrine ($p = 0.5$)}

Setting $p = 0.5$ causes \ps{} to return \ms{}: the plan at the median of
the $T$-plan bisection distribution.
This is the plan that a uniformly random seed would produce with 50\%
probability within the bisection family.

The legal foundation is a \emph{representativeness} norm: the plan should
be typical of what a neutral party running the same algorithm would produce.
A legislature adopting $p = 0.5$ can argue that its plan is not
cherry-picked---it is the plan that an impartial random draw would produce
half the time.

\begin{statute}
\textbf{Representativeness Posture.}
The Districting Administrator shall produce the plan at the median of the
$T$-plan bisection distribution, ranked by normalised edge-cut fraction.
This plan is equally likely to be produced by any party running the
algorithm with a uniformly random seed from the statutory seed space.
Note: this representativeness claim holds within the bisection family
(plans produced by the same SHA-256 seed walk).
It is not a claim that the plan is typical among all valid redistricting
plans---that stronger claim requires ensemble comparison (see G.1).
\end{statute}

The $p = 0.5$ posture defuses the over-optimising objection entirely: the
plan is not an extremum, it is the median.
However, it concedes compactness.
For states where the bisection distribution is already tightly clustered
(B.7 CV $< 2\%$), this concession is minimal.
For states with higher variance (GA CV $= 4.3\%$), it may involve a
measurable compactness reduction.

Crucially, the insensitivity finding confirms that this compactness concession
does not affect partisan outcomes.
The $p = 0.5$ plan produces the same seat counts as the $p = 0.0$ plan.
If the goal is partisan neutrality, both postures achieve it equally.

\subsection{Posture 3: Full Ensemble Representativeness (TargetedSweep)}

A third posture, \ts{}(50th, \recom{}), targets the 50th percentile of the
full \recom{} ensemble~\citep{deford2021} rather than the bisection
distribution.
This posture produces the plan that is ``typical'' in the broadest sense:
typical of all valid redistricting plans, not just those in the bisection
family.

\ts{}(50th, \recom{}) is conceptually appealing but operationally complex.
It requires running a \gc{} \recom{} ensemble (typically 10,000+ steps per
state~\citep{herschlag2020,fifield2020}), computing summary statistics for
each plan, and identifying the plan at the ensemble median.
Unlike \ps{}, it is not a deterministic function of a single statutory
parameter---it depends on the ensemble initialisation and mixing time.

\begin{statute}
\textbf{Full Ensemble Representativeness Posture.}
The Districting Administrator shall produce a plan at the 50th percentile
of the \recom{} Markov-chain ensemble, run for $M$ steps from the
content-derived initialisation.
This plan is typical of all valid redistricting plans consistent with
population balance and contiguity.
\end{statute}

This posture resolves the over-optimising objection most completely but
introduces new vulnerabilities.
The Markov chain must be shown to have mixed~\citep{deford2021}.
The ensemble definition (which constraints are imposed, how plans are scored)
must be specified in statute.
And the resulting plan need not be compact: as G.1 shows, most \recom{}
ensemble plans have higher edge cuts than the \ar{} plan, so a median
\recom{} plan may be substantially less compact.

We regard \ts{}(50th, \recom{}) as a research direction rather than an
immediate statutory option.
The \ps{} framework with $p \in \{0.0, 0.5\}$ covers the legally relevant
range with far simpler implementation.

\subsection{Recommendation: $p = 0.0$ for the DIA}

We recommend that the \dia{} specify $p = 0.0$ (the compactness extremum)
for three reasons.

\textbf{First, legal strength.}
The $p = 0.0$ plan is the most compact consistent with the algorithm.
Under \karcher{}'s logic, this is the strongest possible compactness
posture.
Any challenger who claims the plan is too compact must explain why a less
compact plan would better serve any legitimate redistricting objective.
Given that compactness reduces gerrymandering opportunities and improves
community cohesion, this is a very difficult argument to make.

\textbf{Second, determinism.}
The $p = 0.0$ plan is the global optimum of the bisection family.
It can be certified by running all $T$ seeds and verifying that no seed
achieves a lower $\ECnorm$.
This certification is straightforward and does not depend on any statistical
model.
For $p > 0$, the result depends on the particular $T$-seed sample; for
$p = 0.0$, it is robust to the choice of $T$ (as long as $T \geq T_{\mathrm{stat}}
= 600$, certified in B.16~\citep{b16paper}).

\textbf{Third, partisan neutrality.}
The insensitivity finding (Section~\ref{sec:results}) confirms that
partisan outcomes are the same at $p = 0.0$ as at any other percentile.
The choice of $p = 0.0$ is not a partisan choice; it is a compactness
choice with no partisan consequence.
This neutrality is itself a legal asset: the choice cannot be characterised
as motivated by partisan calculation.

The recommendation is therefore:
\begin{statute}
\textbf{Statutory Recommendation.}
The \dia{} shall specify \ps{}$(0.0, 600)$, equivalent to \cs{} with
$T_{\mathrm{stat}} = 600$.
The legal rationale is the compactness doctrine: the resulting plan is the
most compact consistent with population balance and contiguity within the
\ar{} bisection family.
The statutory record shall state that partisan outcomes are invariant to
the compactness percentile (H.0), so the compactness choice is not a
partisan choice.
\end{statute}
